\section{Introdução}
\label{sec:introducao}

A administração pública brasileira exerce papel central na prestação de serviços essenciais à população. No âmbito municipal, a execução de políticas públicas cotidianas demanda a aquisição de insumos, a contratação de obras e serviços e a mobilização de recursos públicos. Essas atividades estabelecem uma relação contínua entre órgãos contratantes e fornecedores privados, cuja integridade depende, entre outros fatores, da preservação da isonomia e da efetiva competitividade dos procedimentos.

A legislação brasileira estabelece princípios e mecanismos destinados a proteger a competição nas contratações públicas. Destacam-se a Lei nº 14.133/2021, que disciplina as licitações e os contratos administrativos~\cite{brasil2021licitacoes}; a Lei nº 12.529/2011, que estrutura o Sistema Brasileiro de Defesa da Concorrência~\cite{brasil2011concorrencia}; e a Lei nº 12.846/2013, que dispõe sobre a responsabilização administrativa e civil de pessoas jurídicas pela prática de atos contra a Administração Pública~\cite{brasil2013anticorrupcao}. Apesar desse arcabouço, relações societárias, familiares ou cadastrais entre concorrentes formalmente distintos podem justificar verificações adicionais quando coexistem com padrões de participação em certames.

A identificação desses padrões é dificultada pela natureza fragmentada dos registros administrativos. Consultas isoladas a atas de habilitação, julgamento ou cadastros empresariais podem não revelar relações indiretas, como sócios ou controladores compartilhados, vínculos por meio de pessoas jurídicas intermediárias, endereços coincidentes, telefones comuns ou sobreposições onomásticas. Tais coincidências, contudo, não demonstram automaticamente coordenação ilícita: devem ser interpretadas em conjunto com o contexto das disputas e com outros elementos verificáveis.

A Teoria dos Grafos oferece instrumentos para representar essas relações de maneira integrada. Neste trabalho, órgãos públicos, empresas licitantes e sócios ou controladores são modelados como tipos distintos de vértices, enquanto participação em licitações, co-disputa, propriedade societária e vínculos suspeitos são representados em camadas relacionais específicas. Essa modelagem heterogênea multicamada permite examinar simultaneamente relações de mercado e de propriedade, bem como reduzir a rede a subestruturas selecionadas por critérios explícitos.

A motivação prática da pesquisa decorre da disponibilidade de dados relacionais do Sistema Informatizado de Contas dos Municípios (SICOM), mantidos pelo Tribunal de Contas do Estado de Minas Gerais (TCE-MG)~\cite{tce_mg_dados_abertos,tcemg2025in01}. Para o município de Divinópolis/MG, no exercício de 2025, esses registros permitem reconstruir relações entre órgãos e fornecedores, mas os dados societários disponibilizados pelo SICOM não apresentam, por si sós, todos os identificadores necessários à análise de pessoas vinculadas às empresas. Para complementar essa informação, utiliza-se o Quadro de Sócios e Administradores (QSA) disponibilizado por meio da BrasilAPI~\cite{brasilapi2026cnpj}.

O objetivo geral deste trabalho é desenvolver e aplicar um \textit{pipeline} computacional baseado em grafos heterogêneos multicamada para identificar e priorizar, de forma indiciária, pares de empresas concorrentes que apresentem sobreposições societárias, cadastrais ou onomásticas relevantes nas contratações públicas de Divinópolis/MG em 2025. A análise concentra-se na reconstrução da rede licitatória e societária, na atribuição de escores determinísticos aos pares co-disputantes e na redução topológica para identificação de núcleos prioritários. Os resultados não constituem prova automática de fraude ou conluio, mas instrumentos de apoio à seleção de casos para verificação humana.

De modo específico, a investigação contempla: (i) saneamento de documentos fiscais e consolidação dos participantes registrados nas fases de habilitação e julgamento; (ii) construção do grafo multicamada a partir das relações licitatórias do SICOM; (iii) enriquecimento cadastral pelo QSA, com armazenamento local das consultas e aplicação de filtros para reduzir falsos positivos; (iv) cálculo do escore multidimensional exclusivamente para pares de empresas co-disputantes; e (v) redução determinística da rede, identificação de componentes conexas e enumeração de motivos triangulares de compartilhamento societário. O cruzamento financeiro entre certames, empenhos e pagamentos não é apresentado como resultado empírico neste artigo, devido às limitações de vinculação entre módulos identificadas na base; sua integração é tratada como extensão futura.

O restante do artigo está organizado da seguinte forma. A Seção~\ref{sec:referencial} apresenta o referencial jurídico, os estudos sobre redes de co-disputa e relações societárias, e os fundamentos de grafos multicamada e redução topológica. A Seção~\ref{sec:metodologia} descreve as fontes de dados, a construção do grafo, o escore e os procedimentos algorítmicos. A Seção~\ref{sec:resultados} apresenta e discute os resultados obtidos. Por fim, a Seção~\ref{sec:conclusao} sintetiza as contribuições, limitações e possibilidades de continuidade.
